% ─── Capítulo: Pruebas y Demostraciones ───────────────────────────────────────
\chapter{PRUEBAS Y DEMOSTRACIONES DEL COMPILADOR}

\section{Descripción de las Pruebas}

Para validar el funcionamiento integral del compilador TPP se diseñaron programas de prueba que ejercitan cada fase del pipeline. Cada prueba muestra: el código fuente en TPP, el tipo de resultado esperado (correcto o error), y la salida producida por el compilador.

% ─── Macro para placeholder de imagen ─────────────────────────────────────────
\newcommand{\imgplaceholder}[2]{%
    \begin{figure}[H]
    \centering
    \fbox{\begin{minipage}{0.85\textwidth}
        \centering\vspace{1.5cm}
        \textcolor{gray}{\large [Captura de pantalla pendiente]}\\[0.5em]
        \textcolor{gray}{\normalsize #2}
        \vspace{1.5cm}
    \end{minipage}}
    \caption{#1}
    \label{fig:#2}
    \end{figure}
}

% Versión condicional: si el archivo existe, lo muestra; si no, muestra placeholder
\newcommand{\imgtest}[3]{%
    \IfFileExists{img/#1}{%
        \begin{figure}[H]
        \centering
        \includegraphics[width=0.85\textwidth]{img/#1}
        \caption{#2}
        \label{fig:#3}
        \end{figure}
    }{%
        \imgplaceholder{#2}{#3}
    }%
}

% ──────────────────────────────────────────────────────────────────────────────
\section{Prueba 1: Programa Correcto — Operaciones Básicas}

\subsection{Código Fuente}

\begin{lstlisting}[style=tppstyle, caption={Programa correcto: operaciones básicas}]
// Prueba básica: variables, operaciones y salida
entero a = 10;
entero b = 3;
entero resultado = a + b;
imprimir("Suma:");
imprimir(resultado);

resultado = a * b;
imprimir("Producto:");
imprimir(resultado);
\end{lstlisting}

\subsection{Salida Esperada del Compilador}

El compilador procesa el programa sin errores, muestra la tabla de símbolos, el TAC y el ensamblador generado:

\begin{verbatim}
==================== TABLA DE SIMBOLOS ====================
Nombre         Categoria  Tipo    Ambito  Tam  Memoria
a              Variable   entero  0       1    0x00000000
b              Variable   entero  0       1    0x00000004
resultado      Variable   entero  0       1    0x00000008
Total: 3 simbolos | Memoria usada: 12 bytes
==================== OPTIMIZACION ====================
  - Plegado de constantes: 0
==================== CODIGO FINAL (RV32I) ====================
    li t0, 10
    sw t0, 0(sp)
    li t0, 3
    sw t0, 4(sp)
    ...
\end{verbatim}

\imgtest{prueba01sintactico.png}{Ejecución del programa mostrando el árbol sintáctico}{prueba01sintactico}
\imgtest{prueba01semantico.png}{Ejecución del programa mostrando análisis semántico y tabla de símbolos}{prueba01semantico}
\imgtest{prueba01ir.png}{Ejecución del programa mostrando el código intermedio (TAC)}{prueba01ir}
\imgtest{prueba01opt.png}{Ejecución del programa mostrando las optimizaciones aplicadas}{prueba01opt}
\imgtest{prueba01codfin.png}{Ejecución del programa mostrando el código final generado}{prueba01codfin}

% ──────────────────────────────────────────────────────────────────────────────
\section{Prueba 2: Error Léxico}

\subsection{Código Fuente con Error}

\begin{lstlisting}[style=tppstyle, caption={Programa con error léxico: carácter inválido \_}]
// El guión bajo no es válido en TPP
entero mi_variable = 5;   // ERROR: '_' no es válido
imprimir(mi_variable);
\end{lstlisting}

\subsection{Salida del Compilador}

Al encontrar el carácter \texttt{\_}, el lexer reporta el error léxico con línea y columna exactas:

\begin{verbatim}
Error Léxico en la línea 2, columna 11: Carácter inválido '_'
Error Léxico en la línea 2, columna 20: Carácter inválido '_'
\end{verbatim}

El compilador \textbf{no aborta}, sino que reporta todos los errores léxicos encontrados y continúa con los tokens válidos. Sin embargo, los identificadores afectados no serán reconocidos, lo que probablemente dispare errores sintácticos en cascada.

\imgtest{prueba02.png}{Captura de error léxico por carácter inválido}{error_lexico}

% ──────────────────────────────────────────────────────────────────────────────
\section{Prueba 3: Error Sintáctico}

\subsection{Código Fuente con Error}

\begin{lstlisting}[style=tppstyle, caption={Programa con error sintáctico: falta paréntesis}]
// Falta el paréntesis de cierre en el si
entero x = 5;
si (x > 0 {      // ERROR: falta ')'
    imprimir("positivo");
}
\end{lstlisting}

\subsection{Salida del Compilador}

Bison detecta el error sintáctico en la posición exacta. La recuperación de errores (\textit{panic mode}) con el token \texttt{error} permite que el parser continúe analizando el resto del archivo:

\begin{verbatim}
¡Ups! Hay un problema sintáctico en tu código.
Linea 3, Columna 10
Detalle: syntax error, unexpected '{', expecting ')'
=> [Panic Mode] Error en condicion del SI. Recuperando...
\end{verbatim}

\imgtest{prueba03.png}{Captura de error sintáctico por paréntesis faltante}{error_sintactico}

% ──────────────────────────────────────────────────────────────────────────────
\section{Prueba 4: Error Semántico}

\subsection{Código Fuente con Error}

\begin{lstlisting}[style=tppstyle, caption={Programa con error semántico: variable no declarada}]
// Se intenta usar una variable que no fue declarada
entero x = 5;
entero resultado = x + y;  // ERROR: 'y' no declarada
imprimir(resultado);
\end{lstlisting}

\subsection{Salida del Compilador}

El analizador semántico detecta el uso de un identificador que no existe en la tabla de símbolos:

\begin{verbatim}
Error Semantico: Uso de variable no declarada: 'y'
Error Semantico (linea aprox.): Uso de variable no declarada: 'y'
==================== TABLA DE SIMBOLOS ====================
x              Variable   entero  0       1    0x00000000
resultado      Variable   entero  0       1    0x00000004
Total: 2 simbolos | Memoria usada: 8 bytes
[COMPILACION ABORTADA: se encontraron errores semanticos]
\end{verbatim}

\imgtest{prueba04.png}{Captura de error semántico por variable no declarada}{error_semantico}

% ──────────────────────────────────────────────────────────────────────────────
\section{Prueba 5: Programa con Funciones}

\subsection{Código Fuente}

\begin{lstlisting}[style=tppstyle, caption={Programa correcto con declaración y llamada de función}]
fun duplicar(entero n) -> entero {
    retornar n * 2;
}

fun main() {
    entero valor = 7;
    entero doble = duplicar(valor);
    imprimir("Doble de 7:");
    imprimir(doble);
}
\end{lstlisting}

\subsection{Salida Esperada}

\begin{verbatim}
==================== TABLA DE SIMBOLOS ===============================
Nombre               Categoria  Tipo       Ambito   Tam      Memoria   
----------------------------------------------------------------------
.ra_main             Variable   entero     0        1        0x00000008
main                 Funcion    ninguno    0        -        0x00000000
.ra_duplicar         Variable   entero     0        1        0x00000000
duplicar             Funcion    entero     0        -        0x00000000
----------------------------------------------------------------------
Total: 4 simbolos | Memoria usada: 20 bytes (5 words)
\end{verbatim}

\imgtest{prueba05.png}{Ejecución del programa con funciones}{prueba_funciones}

% ──────────────────────────────────────────────────────────────────────────────
\section{Prueba 6: Optimización de Constantes}

\subsection{Código Fuente}

\begin{lstlisting}[style=tppstyle, caption={Programa que activa plegado de constantes y simplificación algebraica}]
entero a = 10 * 2 + 3;   // Debe plegarse a 23
entero b = a + 0;         // Simplificacion algebraica: b = a
entero c = 1 * a;         // Simplificacion algebraica: c = a
imprimir(a);
imprimir(b);
imprimir(c);
\end{lstlisting}

\subsection{Salida del Optimizador}

\begin{verbatim}
==================== OPTIMIZACION DE CODIGO ====================
  Optimizaciones aplicadas:
    - Plegado de constantes:        2  (10*2=20, 20+3=23)
    - Simplificacion algebraica:    2  (a+0->a, 1*a->a)
    - Eliminacion de codigo muerto: 0
  Total: 4 optimizaciones
\end{verbatim}

\imgtest{prueba06.png}{Ejecución mostrando optimizaciones aplicadas}{prueba_optimizacion}

% ──────────────────────────────────────────────────────────────────────────────
\section{Prueba 7: Programa con Bucles y Arreglos}

\subsection{Código Fuente}

\begin{lstlisting}[style=tppstyle, caption={Programa con arreglo y bucle para}]
entero arr[5];
entero i;
para (entero k = 0; k < 5; k = k + 1) {
    arr[k] = k * 2;
}
para (entero j = 0; j < 5; j = j + 1) {
    imprimir(arr[j]);
}
\end{lstlisting}

\subsection{Salida Esperada}

El programa imprime: 0, 2, 4, 6, 8 (los primeros 5 múltiplos de 2).

\imgtest{prueba07logism.png}{Ejecución del programa con arreglos y bucles para}{prueba_arreglos}

% ──────────────────────────────────────────────────────────────────────────────
\section{Prueba 8: Prueba Integral — Backend}

\subsection{Código Fuente}

\begin{lstlisting}[style=tppstyle, caption={Prueba backend: sumar elementos de arreglo}]
fun sumarArreglo(entero n) -> entero {
    entero suma = 0;
    entero arr[5];
    arr[0] = 1;
    arr[1] = 2;
    arr[2] = 3;
    arr[3] = 4;
    arr[4] = 5;
    para (entero i = 0; i < n; i = i + 1) {
        suma = suma + arr[i];
    }
    retornar suma;
}

fun main() {
    entero total = sumarArreglo(5);
    imprimir("Suma 1..5:");
    imprimir(total);
}
\end{lstlisting}

\subsection{Resultado en Logisim}

Con el archivo \texttt{output.hex} cargado en la RAM del procesador Logisim, la ejecución imprime \texttt{15} por el TTY del circuito.

\imgtest{prueba08logism.png}{Ejecución del programa en Logisim Evolution — salida TTY}{prueba_logisim}
